\section{Johdanto}

Kielimallit ovat luonnollisen kielen käsittelyn malleja, joiden avulla
voidaan kuvata tekstin tilastollista rakennetta tai tuottaa uutta tekstiä.
Autoregressiivisessa kielimallinnuksessa malli ennustaa seuraavaa sanaa
(tai tokenia eli tekstin käsittely-yksikköä) aiemman kontekstin perusteella
\cite{jurafsky2023speech, bengio2003neural}. Tekstiä voidaan generoida
valitsemalla mallin ennustamasta jakaumasta seuraava token, liittämällä se
kontekstiin ja toistamalla menettelyä päivitetyllä kontekstilla. Tämä periaate yhdistää klassisia
n-gram-malleja \cite{shannon1948mathematical} ja nykyaikaisia transformer pohjaisia kielimalleja \cite{vaswani2017attention}, vaikka niiden
tavat esittää kontekstia ja muodostaa ennustejakauma eroavat huomattavasti.

Kielimallien kehityksessä keskeinen muutos on ollut siirtyminen
diskreeteistä tekstiesityksistä opittuihin jatkuviin esityksiin \cite{mikolov2013efficient} sekä
esiintymismääriin perustuvista malleista neuroverkkopohjaiseen
mallinnukseen \cite{werbos1974beyond, rumelhart1986learning}. Upotusvektoreiden avulla tokenit esitetään vektoriavaruudessa, jossa niiden väliset suhteet voidaan oppia aineistosta. Samalla neuroverkon parametreja optimoidaan tuottamaan opetusaineistoa vastaavia jakaumia. Nykyaikaisissa
dekooderipohjaisissa Transformer-malleissa kontekstin käsittely perustuu
itsehuomioon \cite{vaswani2017attention}, jonka avulla tokenien esityksiä muodostetaan suhteessa
muihin kontekstin tokeneihin. Kielimallien toiminnan ymmärtäminen edellyttää
siten tekstiesitysten, todennäköisyysmallinnuksen, optimoinnin ja
neuroverkkoarkkitehtuurin tarkastelua yhtenä kokonaisuutena.

Suurten kielimallien sovellukset eivät rajoitu vapaamuotoisen tekstin
tuottamiseen. Kielimalleja voidaan liittää järjestelmiin, joissa osa niiden
tuottamista tokeneista tulkitaan työkalukutsuina
\cite{schick2023toolformer, yao2023react}. Tällöin kielimallia ympäröivä
ohjelmisto voi suorittaa laskentaa, tehdä verkkohakuja tai käsitellä
tietokantakyselyitä. Yksi käytännössä merkittävä sovelluskohde on
luonnollisen kielen muuntaminen SQL-kyselyiksi eli Text-to-SQL
\cite{liu2024survey, deng2022recent}. Sen tavoitteena on mahdollistaa
tietokantojen hyödyntäminen ilman, että käyttäjän tarvitsee itse kirjoittaa
SQL-kyselyitä tai tuntea tietokannan rakennetta yksityiskohtaisesti.

Text-to-SQL-järjestelmän hyödyllisyys riippuu kuitenkin siitä, kuinka
luotettavasti se kykenee yhdistämään käyttäjän tiedontarpeen tietokannan
rakenteeseen ja muodostamaan tarkoitusta vastaavan kyselyn
\cite{yu2018spider, li2024bird}. Syntaktisesti kelvollinen SQL-kysely ei
välttämättä vastaa käyttäjän tarkoittamaa kysymystä. Luotettavuuden kannalta
on myös olennaista, säilyykö järjestelmän tulkinta johdonmukaisena saman
kysymyksen toistoissa ja silloin, kun kysymyksen esitystapaa muutetaan
sen merkitystä muuttamatta. Nämä kysymykset motivoivat kielimallipohjaisten
järjestelmien empiiristä arviointia.

Tämän tutkielman tavoitteena on jäsentää autoregressiivisen
kielimallinnuksen matemaattisia perusteita ja tarkastella
kielimallipohjaisen Text-to-SQL-järjestelmän toimintaa käytännön
tapaustutkimuksessa. Teoreettisessa osassa rakennetaan yhteys
diskreettien tekstiesitysten ja klassisten n-gram-mallien, jatkuvien
upotusvektorien, neuroverkkopohjaisen ennustamisen sekä
dekooderipohjaisen Transformer-arkkitehtuurin välille. Lisäksi
tarkastellaan gradienttimenetelmää ja adaptiivisia optimointimenetelmiä
kielimallien parametrien oppimisen perustana. Tavoitteena ei ole
kattava katsaus kaikkiin kielimalliarkkitehtuureihin, vaan asteittain
etenevä esitys niistä käsitteistä ja menetelmistä, joihin tarkasteltava
autoregressiivinen kielimallinnus perustuu.

Tutkielman empiirisessä osassa tarkastellaan Snowflaken Cortex
Analyst -järjes-\\telmää \cite{CortexAnalyst}. Se muodostaa käyttäjän
luonnollisen kielen kysymyksistä SQL-kyselyitä hyödyntäen tietokannan
rakennetta ja käsitteitä kuvaavaa semanttista näkymää
\cite{snowflake_cortex_semantic_model}. Tapaustutkimuksen aineistona on
puhelinkeskusdataa sisältävä tietokanta, joka koostuu yhdestä
keskusteludataa sisältävästä päätaulusta ja kolmesta sitä täydentävästä
taulusta. Semanttinen näkymä kuvaa aineiston tauluja, kenttiä ja niiden
merkityksiä ja toimii siten välittävänä kuvauksena luonnollisen kielen
kysymysten ja tietokannan välillä \cite{CortexAnalyst}.

Teoreettinen tarkastelu tarjoaa käsitteellisen perustan
kielimallipohjaisen kyselymuodostuksen ymmärtämiselle, mutta sitä ei
tule tulkita Cortex Analystin sisäisen toteutuksen kuvaukseksi.
Kaupallisten kielimallijärjestelmien tarkat toteutustiedot eivät
yleensä ole kokonaisuudessaan julkisia, joten tapaustutkimuksessa
järjestelmää arvioidaan sen tuottamien SQL-kyselyiden perusteella.
Tarkastelun kohteena on näin ollen järjestelmän havaittava toiminta,
ei yksittäisten mallikomponenttien vaikutusten eristäminen.

Arviointi rajataan tuotettuihin SQL-kyselyihin eikä niiden suorittamisen
palauttamiin tulosjoukkoihin. Rajaus on perusteltu aineiston
luottamuksellisuudella sekä sillä, että virheellisen kyselyn vaikutus
tulosjoukkoon riippuu myös tietokannan sisällöstä. Toisaalta erilaiset
SQL-kyselyt voivat olla merkitykseltään yhtäpitäviä, joten kyselyiden
välinen vaihtelu ei sellaisenaan osoita virhettä. Tarkastelussa
erotetaan siksi kyselyiden vaihtelu ja niiden vastaavuus käyttäjän
tiedontarpeeseen. Tulokset kuvaavat järjestelmän toimintaa valitussa
tutkimusasetelmassa, eivät sen yleistä suorituskykyä kaikissa
Text-to-SQL-tehtävissä.

Tutkielma jakautuu teoriaosaan ja empiiriseen osaan. Luvuissa \ref{sec:discrete}-\ref{sec:transformer} rakennetaan kielimallien matemaattinen ja arkkitehtuurinen perusta diskreeteistä esityksistä moderniin Transformer-malliin ja optimointiin. Luvussa \ref{sec:methodology} esitellään Cortex Analystiin pohjautuva tutkimusasetelma, ja luvuissa \ref{sec:experiments}-\ref{sec:results} kuvataan kokeet, analysoidaan tulokset sekä pohditaan työn rajoituksia ja johtopäätöksiä.
